\documentclass[11pt]{article}

\usepackage[margin=0.85in]{geometry}
\usepackage{amsmath}
\usepackage{graphicx}
\usepackage{hyperref}
\usepackage{cite}
\usepackage{float}

\title{Constraint-Gated Waveform Design for Low-Excitation Ion Transport}
\author{Dan Hawkley}
\date{2026}

\begin{document}
\maketitle

\begin{abstract}
We present a minimal, reproducible pipeline for designing trapped-ion
transport waveforms and directly evaluating their dynamical consequences.
Using a one-dimensional model, we show that smooth minimum-jerk transport
trajectories suppress motional excitation by multiple orders of magnitude
compared to linear interpolation, and we quantify clear tradeoffs between
transport speed, waveform smoothness, and residual excitation. The
framework is intentionally simple so that these relationships are
transparent and directly inspectable, providing a practical bridge
between waveform design and physical performance.
\end{abstract}

\section{Introduction}

Segmented trapped-ion architectures require reliable ion transport between
trap zones. Transport waveforms must move ions quickly while limiting
residual motional excitation. This work develops a small, reproducible
simulation pipeline that directly connects waveform design to excitation
outcomes.

This work focuses on making these tradeoffs directly visible through a
minimal computational model, rather than optimizing a specific hardware
implementation.

\section{Related Work}

Ion transport in segmented traps has been studied extensively in both
experimental and theoretical contexts. Early work demonstrated reliable
transport and splitting operations in RF traps~\cite{rowe2002transport},
while later experiments achieved fast transport with low excitation using
optimized control protocols~\cite{bowler2012coherent}.

From a theoretical perspective, smooth trajectory design has been shown
to suppress excitation during transport~\cite{torrontegui2011fast},
motivating the use of minimum-jerk and related interpolation schemes.

The present work complements this literature by providing a minimal,
reproducible pipeline in which these effects are directly visible and
quantifiable.

\section{Pipeline}

The computational pipeline is:
\[
\text{electrode basis}
\;\rightarrow\;
\text{potential well}
\;\rightarrow\;
\text{transport path}
\;\rightarrow\;
\text{voltage waveform}
\;\rightarrow\;
\text{ion motion}
\;\rightarrow\;
\text{excitation metric}.
\]

\section{Transport Model}

Ion motion is modeled as a particle in a moving harmonic well:
\[
m\ddot{x}(t) = -m\omega^2[x(t)-x_c(t)],
\]
where \(x_c(t)\) is the waveform-controlled well center.

\section{Waveform Design}

Smooth transport paths are generated using minimum-jerk interpolation:
\[
s(u)=10u^3-15u^4+6u^5.
\]
This ensures zero velocity and acceleration at endpoints, reducing
excitation.

\section{Results}

\subsection{Transport Trajectory}

The ion trajectory closely follows the target transport path under
minimum-jerk control, indicating low excitation during motion.

\begin{figure}[H]
\centering
\includegraphics[width=0.78\textwidth]{../figures/04_motion_trace.png}
\caption{Ion trajectory closely follows the target transport path under minimum-jerk control.}
\end{figure}

\subsection{Path Comparison}

Minimum-jerk transport significantly suppresses oscillatory residual
motion compared to linear interpolation, which strongly excites the trap.

\begin{figure}[H]
\centering
\includegraphics[width=0.78\textwidth]{../figures/05_excitation_residual_time_series.png}
\caption{Minimum-jerk transport suppresses oscillatory excitation relative to linear interpolation.}
\end{figure}

\subsection{Excitation Scaling}

Residual excitation decreases systematically with increasing transport
duration, demonstrating a clear tradeoff between speed and excitation.

\begin{figure}[H]
\centering
\includegraphics[width=0.78\textwidth]{../figures/06_excitation_vs_duration.png}
\caption{Residual excitation decreases systematically with increasing transport duration.}
\end{figure}

\subsection{Frequency-Domain Analysis}

Frequency-domain analysis shows that linear transport excites the trap
near the secular frequency, while minimum-jerk trajectories suppress
resonant excitation.

\begin{figure}[H]
\centering
\includegraphics[width=0.78\textwidth]{../figures/05_excitation_fft_spectrum.png}
\caption{Residual-motion spectrum showing suppression of excitation at the trap secular frequency.}
\end{figure}

\clearpage

\section{Discussion}

The key result of this work is that excitation behavior emerges clearly
from trajectory smoothness alone, without requiring complex control
optimization.

These results suggest that smooth trajectory design provides a simple
and effective constraint on excitation, offering a practical design
principle for ion transport systems.

Smooth transport paths suppress resonant excitation by minimizing
dynamical forcing. The results demonstrate:

\begin{itemize}
\item Orders-of-magnitude reduction in excitation for minimum-jerk paths
\item Clear scaling laws linking transport speed and excitation
\item Practical control constraints arising from waveform smoothness
\end{itemize}

\section{Conclusion}

We presented a compact, interpretable pipeline for ion transport waveform
design. The framework directly connects control synthesis to physical
outcomes and highlights fundamental tradeoffs between speed and
excitation. This approach provides a foundation for more advanced optimal
control and experimental implementation.

\section*{Code Availability}

All code and figures are generated from a small, reproducible pipeline
available at:

\href{https://github.com/thinkthoughts/ion-transport-waveform-pipeline}
{github.com/thinkthoughts/ion-transport-waveform-pipeline}

\bibliographystyle{plain}
\bibliography{references}

\end{document}
